\chapter{Importancia de la protección radiológica}

El uso de la radiactividad en la ingeniería es uno de los aspectos más difundidos de la tecnología nuclear, puesto que forma parte del pequeño número de técnicas de control de calidad mediante pruebas no destructivas.

Pero así como son muchas las ventajas que nos ofrece, también se corren muchos riesgos si no sabemos emplearla bajo ciertos procedimientos que garanticen nuestra protección contra la radiación.

Estar conscientes de que el programa de protección radiológica está hecho para asegurar nuestra seguridad es dar el primer paso en la comprensión de este curso.

El objetivo principal de la protección radiológica es lograr que las operaciones controladas por la dirección de un centro se desarrollen sin riesgo para la salud y la seguridad de las personas que estén dentro o fuera del área, así como la protección del medio ambiente, controlando los riesgos que se derivan de las actividades que, por las características de ciertos materiales o equipos que se utilizan, puedan implicar la exposición a radiaciones ionizantes. Para conseguirlo, hay que aplicar los reglamentos que estén en vigor y todas las demás normas pertinentes, y vigilar debidamente las operaciones para comprobar la eficacia de las medidas de protección adoptadas.

Las responsabilidades que se adquieren al ser nombrado como personal ocupacionalmente expuesto son:

\begin{itemize}
  \item Conocer y aplicar correctamente los principios básicos de seguridad radiológica.
  \item Conocer el uso correcto de fuentes de radiación ionizante, equipo detector y accesorios.
  \item Usar durante la jornada de trabajo el vestuario, equipo, accesorios y dispositivos que se requieran, de acuerdo con lo establecido en el manual de seguridad radiológica y lo dispuesto por el encargado de seguridad radiológica.
  \item Informar al encargado de seguridad radiológica sobre cualquier situación de alto riesgo, incidente o accidente radiológico.
\end{itemize}
